\documentclass[a4paper]{article}
\usepackage{amsmath}
\usepackage{amsfonts}
\usepackage{amssymb}

\begin{document}

\noindent \large Auteur: Dina Haajou \\
\noindent \large Studierichting: Informatica \\
\noindent \large Oefeningenreeks 6A nr. 6.3 \\

\medskip
\normalsize

$x_1 + x_4 = 16$ en $x_3 + x_5 = 28$ \\

$x_n = x_1 + (n-1)v$ \\

Voor $x_4$: \\

$x_4 = x_1 + 3v \Rightarrow x_1 + x_4 = x_1 + (x_1 + 3v) = 2x_1 + 3v = 16$ \\

Voor $x_3$ en $x_5$: \\

$x_3 = x_1 + 2v$, $x_5 = x_1 + 4v$ \\

$x_3 + x_5 = x_1 + 2v + x_1 + 4v = 2x_1 + 6v = 28$ \\

Neem het verschil van de vergelijkingen: \\

$(2x_1 + 6v) - (2x_1 + 3v) = 28 - 16 \Rightarrow 3v = 12 \Rightarrow v = 4$ \\

Vul $v = 4$ in: \\

$2x_1 + 3 \cdot 4 = 16 \Rightarrow 2x_1 = 4 \Rightarrow x_1 = 2$ \\


Hieruit volgt: \\

$x_1 = 2$ \\

$x_2 = x_1 + v = 6$ \\

$x_3 = x_1 + 2v = 10$ \\

\begin{thebibliography}{9}
% \bibitem{}
\end{thebibliography}

\end{document}
